\documentclass[12pt]{article}
\usepackage{amsmath,amsfonts,fullpage,amsthm,graphicx}

\newtheorem*{theorem}{Theorem}
\newtheorem*{fact}{Fact}
\newtheorem*{goal}{Goal}
\newtheorem*{idea}{Idea}
\newtheorem*{defn}{Definition}
\newtheorem*{ex}{Example}

%Style section
%\setlength{\textheight}{10in}
%\setlength{\textwidth}{7in}
%\hoffset=-0.75in
 \pagestyle{plain}
% \setlength{\topmargin}{0in}
% \setlength{\headsep}{0in}
% \setlength{\oddsidemargin}{0in}
% \setlength{\evensidemargin}{0in}

\begin{document}


\pagestyle{empty} %use this to get rid of page numbers

\noindent
{Math 166 \hskip 1 truein  Name:\ \hrulefill}

\noindent
%\vskip 0.3truein
%\bigskip
%\noindent
{New series from old}
\noindent \hspace{2.5in} Section Number: \hrulefill
\bigskip


\begin{enumerate}
%\setcounter{enumi}{-1}
\item Find the \emph{radius of convergence} and \emph{interval of convergence} of the power series  
\[\displaystyle\sum_{n=0}^{\infty} (-1)^n \displaystyle\frac{n}{4^n}(x+1)^{2n}.\]


\vfill


\item Find the {radius (and interval) of convergence} for each of the following power series:  $\displaystyle\sum_{n=0}^{\infty} \displaystyle\frac{x^{n}}{n!}$ and  $\displaystyle\sum_{n=0}^{\infty} n!x^{n}$.

\vfill

%\iffalse
\pagebreak
\item A series of the form $\sum_{n = 0}^\infty a_n x^n$ is known to converge at $x = 4$.  For each of the following values, determine whether the series must converge, must diverge, or it is not possible to determine.  Justify your answer.
\begin{enumerate}
\item $x = -3$.
\vspace{1 in}
\item $x = -4$.
\vspace{1 in}
\item $x = 5$.
\vspace{1 in}
\end{enumerate}


%\item Use the power series $\displaystyle\frac{1}{1-x}=\displaystyle\sum_{n=0}^{\infty}x^n$ to find a power series for $\displaystyle\frac{3}{4+7x}$ centered at $0$. Hint: factor the denominator as $4(1-(-\frac{7}{4}x))$. What is the interval of convergence?

%\vspace{2.2in}
\item Find the power series for $\arctan x$.  (Hint: what is the derivative of $\arctan x$?  Express this in the form $\displaystyle \sum_{n = 0}^\infty ar^n = \frac{a}{1 - r}$.)
\vfill
\vfill
\vfill


\item Use the answer to your previous problem to find a series that converges to $\pi$.  (Hint: what is $\arctan 1$?)
\vfill

%What is the radius of convergence? (No ratio test needed - use the information in the table.)
%\fi

\end{enumerate}
\end{document}
